\documentclass[11pt]{article}
\usepackage{preamble}
\renewcommand{\answer}[1]{}

\begin{document}

\ladderheading{AP Statistics \textperiodcentered\ Topic 4.3 (CED 2.4) \textperiodcentered\ B: 2026-10-22 \textperiodcentered\ E: 2026-11-06}{Probability or Estimate?}

\focusbanner{TODAY'S PROBLEM}

Suppose a game uses a spinner with eight equal sectors numbered 1 through 8. A player wins when the spinner lands on 1, 2, or 3. The dotplot shows the results of 40 observed plays.\par\smallskip \textbf{Dani:} ``The probability of winning is $3/8=0.375$ because three of the eight equally likely results win.''\par \textbf{Micah:} ``The probability is $18/40=0.45$ because 18 of the 40 observed plays were wins.''\par
\begin{center}
\begin{tikzpicture}
\begin{axis}[
  width=0.77\linewidth,
  height=1.44in,
  ymin=0, ymax=9, ytick=\empty, axis y line=none, axis x line=bottom,
  xlabel={Spinner result (40 plays)}, tick label style={font=\small}, label style={font=\small}
]
\addplot[only marks, mark=*, mark size=2.2pt, color=dayblue] coordinates {
  (1,1)
  (1,2)
  (1,3)
  (1,4)
  (1,5)
  (1,6)
  (1,7)
  (1,8)
  (2,1)
  (2,2)
  (2,3)
  (2,4)
  (2,5)
  (2,6)
  (3,1)
  (3,2)
  (3,3)
  (3,4)
  (4,1)
  (4,2)
  (4,3)
  (4,4)
  (4,5)
  (5,1)
  (5,2)
  (5,3)
  (6,1)
  (6,2)
  (6,3)
  (6,4)
  (6,5)
  (6,6)
  (7,1)
  (7,2)
  (7,3)
  (7,4)
  (8,1)
  (8,2)
  (8,3)
  (8,4)
};
\end{axis}
\end{tikzpicture}
\end{center}
\begin{firsttakebox}{FIRST TAKE --- before the video (5 min)}
Which number would you use to describe the chance of winning the next play: $0.375$ or $0.45$? Give one reason.
\workspace{0.9in}
\end{firsttakebox}

\begin{callout}{\IconBook\ MODEL, DATA, AND THE LONG RUN (keep on the page)}{calloutgreen}
For equally likely outcomes, $P(E)=\dfrac{\text{outcomes in }E}{\text{outcomes in the sample space}}$.
A \textbf{theoretical probability} comes from a probability model; an \textbf{empirical probability}
is the observed relative frequency $\dfrac{\text{times }E\text{ occurred}}{\text{number of trials}}$ and estimates it.
The \textbf{law of large numbers} says that, as the number of independent trials made under the same
conditions grows, the observed relative frequency tends to approach the true probability. It does not
guarantee an exact match in a fixed number of trials. A fair model requires its stated outcomes to be equally likely.
\end{callout}

\newpage
\textbf{Finish after the video:}\par\smallskip

\dokbadge{1}~\textbf{(a)} List the sample space and identify the outcomes in the event `win.'\par
\workspace{0.7in}

\dokbadge{2}~\textbf{(b)} Use the dotplot to verify the observed relative frequency of a win. Compute the theoretical probability under the equal-sector model, then give the difference in percentage points.\par
\workspace{1.3in}

\dokbadge{3}~\textbf{(c)} Under the stated equal-sector model, decide which number is the theoretical probability and which is an empirical estimate. Cite both the sample-space feature and the 40-play result, use the law of large numbers to explain why the numbers can differ, and state two features of the trial process you would check before trusting the observed proportion.\par
\workspace{2.0in}

\begin{sentenceframebox}\raggedright\textbf{Frame:} The value \blank[0.7in] is the theoretical probability because \blank[2.5in], while \blank[0.7in] is an estimate because \blank[2.1in].\end{sentenceframebox}

\begin{sentenceframebox}\raggedright\textbf{Frame:} The values can differ because \blank[2.6in]; before trusting the estimate, I would check \blank[2.5in].\end{sentenceframebox}

\begin{summaryexitbox}{EXIT --- turn this in}
One thing the video changed about my first take: \hrulefill\par
\workspace{0.4in}
\end{summaryexitbox}

\end{document}
